% architecture.tex — three-tier architecture overview for ComplexGitSync

\chapter{Architecture Overview}
\label{chap:architecture}

This chapter describes how \cgspkg{} is structured around three explicit tiers
and explains the interaction between its classes. Related domain classes are
grouped in modules according to the responsibility boundaries below.

% -----------------------------------------------------------------------
\section{Three-Tier Architecture}
\label{sec:three-tier}

\cgspkg{} is organised into three tiers.  Dependencies only flow
\textbf{downward}: the Client calls Actions; Actions read and mutate Core Data;
Core Data has no upward dependency.

Figure~\ref{fig:three-tier} shows the overall layout.

\begin{figure}[ht]
\centering
\caption{Three-tier architecture: Client/API exposes state-gated methods;
  Actions operate on Core Data; Core Data separates reference identities from
  the \texttt{WorkingGitTree} parent--leaf runtime graph.}
\label{fig:three-tier}
% three_tier.tex — TikZ figure showing the three-tier architecture
% Included by docs/architecture.tex (Figure three-tier).

\begin{tikzpicture}[
  tier/.style={
    rectangle, rounded corners=6pt,
    minimum width=11cm, minimum height=1.4cm,
    draw, thick, align=center,
    font=\sffamily\small,
  },
  tierlabel/.style={
    font=\sffamily\bfseries\small,
    rotate=0,
  },
  downarrow/.style={-{Stealth[length=8pt]}, thick, gray},
]

% Tier 3 — Client / API  (top)
\node[tier, fill=blue!12, draw=blue!40] (api) at (0,0) {
  \textbf{Tier 3 — Client / API}\\[2pt]
  \texttt{ComplexGitSyncClient} \ $\cdot$ \ \texttt{Orchestre} \ $\cdot$ \ \texttt{CLI}\\
  \textit{Exposes class methods; gates every call on \texttt{TreeLifecycleState}}
};

% arrow down
\draw[downarrow] (api.south) -- ++(0,-0.5) node[midway, right, font=\scriptsize\itshape]{calls (state-gated)};

% Tier 2 — Actions (middle)
\node[tier, fill=green!10, draw=green!40] (actions) at (0,-2.3) {
  \textbf{Tier 2 — Actions}\\[2pt]
  \texttt{clone} \ $\cdot$ \ \texttt{checkout} \ $\cdot$ \ \texttt{commit} \ $\cdot$
  \ \texttt{push} \ $\cdot$ \ \texttt{tag} \ $\cdot$ \ \texttt{freeze\_release}\\
  \textit{Each action is available only from valid \texttt{TreeLifecycleState} values}
};

% arrow down
\draw[downarrow] (actions.south) -- ++(0,-0.5) node[midway, right, font=\scriptsize\itshape]{reads / mutates};

% Tier 1 — Core Data (bottom)
\node[tier, fill=orange!10, draw=orange!40] (core) at (0,-4.6) {
  \textbf{Tier 1 — Core Data}\\[2pt]
  \texttt{GitTree} \ $\cdot$ \ \texttt{GitRepo} \ $\cdot$ \ \texttt{DependencyTreeRegistry}
  \ $\cdot$ \ \texttt{RepoRegistryEntry}\\
  \textit{Centred on \texttt{GitTree} and its \texttt{GitRepo} parent--leaf graph}
};

\end{tikzpicture}

\end{figure}

% -----------------------------------------------------------------------
\section{Architectural Positioning}
\label{sec:architectural-positioning}

\cgspkg{} is positioned as a deterministic synchronisation layer \emph{on top of
Git}, not as a replacement for Git itself. It combines two distinct graph
scopes:

\begin{itemize}
  \item \textbf{Git DAG}: per-repository commit history and native remote
    semantics.
  \item \textbf{GitTree DAG}: workspace-level dependency graph coordinating
    parent/root/leaf propagation and ordering.
\end{itemize}

The system therefore differs from:
\begin{itemize}
  \item plain Git usage (single repository scope),
  \item monorepo layouts (single storage boundary),
  \item submodule-only management (linking primitive without lifecycle
    orchestration),
  \item generic distributed sync engines (not constrained by Git-native
    `.gts`/`.lgr` deterministic contracts).
\end{itemize}

Figure~\ref{fig:positioning-matrix} summarizes this positioning.

\begin{figure}[ht]
\centering
\caption{Architectural positioning matrix across Git scope and workspace
  coordination guarantees.}
\label{fig:positioning-matrix}
% positioning_matrix.tex — table-style matrix for architectural positioning.
% Included by docs/Text/architecture.tex.

\begin{tabular}{p{3.2cm}p{4.4cm}p{6.6cm}}
\textbf{Approach} & \textbf{Primary scope} & \textbf{What it does not guarantee} \\
\hline
Git (single repo) & Commit DAG per repository & Multi-repository propagation and coordinated lifecycle gating \\
Monorepo & One repository, one commit DAG & Independent repository identities and per-repo remote ownership \\
Submodule management & Parent repo references child revisions & Deterministic tree-wide mutation workflow and snapshot contracts \\
ComplexGitSync & Git DAG + GitTree DAG & Not a credential broker or remote policy authority \\
\hline
\end{tabular}

\end{figure}

\texttt{freeze} materializes deterministic checkpoints as
\texttt{.gts} snapshots (schema + SHA-256 hash), while \texttt{.lgr}
assigns stable local ids and stores append-only ledger events for replayable
local evolution.

Cross-declared nested repositories may initially resemble a local tangle at
declaration time; runtime expansion normalizes this into a DAG via
\texttt{fix\_circularities}.

% -----------------------------------------------------------------------
\section{Tier 1 — Core Data}
\label{sec:tier-core}

The reference model is centred on \textbf{\texttt{GitTree}}, which owns
canonical \textbf{\texttt{GitRepo}} identities. Its runtime counterpart,
\textbf{\texttt{WorkingGitTree}}, adds the parent--leaf graph of mutable
\textbf{\texttt{WorkingRepo}} nodes. A project consists of one \emph{root}
repository that may nest \emph{parent} repositories and terminal \emph{leaf}
repositories.

Figure~\ref{fig:gittree-nodes} shows a representative project tree.

\begin{figure}[ht]
\centering
\caption{\texttt{WorkingGitTree} is a parent--leaf graph of
  \texttt{WorkingRepo} nodes.
  Root and parent nodes own nested \texttt{.cgs} files discovered during
  \texttt{clone} or \texttt{discover\_nested\_configs}.}
\label{fig:gittree-nodes}
% gittree_nodes.tex — TikZ figure showing GitRepo parent-leaf tree structure
% Included by docs/architecture.tex (Figure gittree-nodes).

\begin{tikzpicture}[
  reponode/.style={
    rectangle, rounded corners=3pt,
    minimum width=3.4cm, minimum height=0.9cm,
    draw, thick, align=center,
    font=\ttfamily\footnotesize,
    inner sep=4pt,
  },
  rootnode/.style={reponode, fill=blue!15, draw=blue!50},
  parentnode/.style={reponode, fill=green!12, draw=green!50},
  leafnode/.style={reponode, fill=orange!12, draw=orange!40},
  treearrow/.style={-{Stealth[length=5pt]}, thick, gray!70},
  label/.style={font=\sffamily\scriptsize\itshape, text=gray!70},
]

% Root
\node[rootnode] (root) at (0,0) {
  \textbf{root} (project)\\NodeType = ROOT
};

% Two direct children of root
\node[parentnode] (parent) at (-3,-2) {
  \textbf{child-repo}\\NodeType = PARENT
};
\node[leafnode] (lib) at (3,-2) {
  \textbf{lib}\\NodeType = LEAF
};

% Children of parent
\node[leafnode] (sub1) at (-4.5,-4) {
  \textbf{docs}\\NodeType = LEAF
};
\node[leafnode] (sub2) at (-1.5,-4) {
  \textbf{fixtures}\\NodeType = LEAF
};

% Tree edges
\draw[treearrow] (root.south) -- (parent.north);
\draw[treearrow] (root.south) -- (lib.north);
\draw[treearrow] (parent.south) -- (sub1.north);
\draw[treearrow] (parent.south) -- (sub2.north);

% Annotations
\node[label, right=0.3cm of root] {owns nested .cgs (root)};
\node[label, right=0.2cm of parent] {owns nested .cgs (parent)};
\node[label, right=0.2cm of lib] {terminal — no nested .cgs};
\node[label, right=0.2cm of sub1] {nested\_config = "disabled"};
\node[label, right=0.2cm of sub2] {nested\_config = "disabled"};

% Legend
\node[rootnode,  minimum width=2.0cm, minimum height=0.6cm]
    (leg1) at (-4.5,-5.5) {\textbf{ROOT}};
\node[parentnode, minimum width=2.0cm, minimum height=0.6cm]
    (leg2) at (-1.5,-5.5) {\textbf{PARENT}};
\node[leafnode,  minimum width=2.0cm, minimum height=0.6cm]
    (leg3) at ( 1.5,-5.5) {\textbf{LEAF}};

\node[font=\sffamily\scriptsize, above=0.1cm of leg1] {Root node};
\node[font=\sffamily\scriptsize, above=0.1cm of leg2] {Parent node};
\node[font=\sffamily\scriptsize, above=0.1cm of leg3] {Leaf node};

\end{tikzpicture}

\end{figure}

Registry entries in \texttt{WorkingGitTree} are indexed by a
colon-separated repo ID derived from their position in the tree (e.g.\
\texttt{root:deps/child-repo:docs}).
Cross-referenced parents therefore create a tangle-like declaration layer, but
the expanded runtime graph remains a DAG after the cycle-breaking engine has
run.

\subsection*{Cycle Breaking Engine}

\texttt{fix\_circularities} is a two-phase engine that converts a potentially
cyclic dependency graph into a valid DAG:

\begin{enumerate}
  \item \textbf{SCC detection (Tarjan's algorithm).}  A directed dependency
    graph is built from the registry (edges: parent absolute path $\to$ child
    absolute path).  \texttt{find\_strongly\_connected\_components} identifies
    groups of paths that form a cycle.  For each non-trivial SCC, an
    \emph{Anchor} path is selected using three heuristics: (1)~most external
    incoming edges, (2)~closest to project root, (3)~smallest SHA-256 hash.
    Back-edge entries — registry entries resolving to the anchor's path but
    parented inside a non-anchor SCC node — are flagged
    \texttt{is\_external\_reference\,=\,True} and removed.
  \item \textbf{Hash-compatibility deduplication.}  Residual path-duplicate
    entries are removed when their synchronisation state is compatible with
    the canonical entry.
\end{enumerate}

The clone engine uses a \textbf{sync stack} (set of absolute paths already
entered into the clone pipeline) to exclude entries whose path is already
being processed, providing defence-in-depth against any residual back-references
that escape the registry cleanup phase.

\texttt{topological\_sort(registry)} returns entries in parent-first order
(Kahn's BFS algorithm) for safe sequential clone/pull operations.  Runtime
\texttt{pull} uses the same three-level direction: \texttt{ROOT -> PARENT ->
LEAF}; every repository --- root, parent, and leaf alike --- is a plain
independent clone and is pulled in that same parent-first order, rather than
through any submodule linkage.

Figure~\ref{fig:core-model} shows the class overview across all three
tiers, \texttt{GitTree}/\texttt{GitRepo} among them; the full current
module-by-module Ring-level dependency graph lives in
\texttt{docs/DevGuide/architecture.md}.

% -----------------------------------------------------------------------
\section{Tier 2 — Actions}
\label{sec:tier-actions}

Actions are operations that read or mutate the Core tier.  \textbf{Every
action is gated by the current \texttt{TreeLifecycleState}}: actions that
require a \texttt{READY} tree will refuse to execute and log a gating-refusal
event if the tree is not ready.

Figure~\ref{fig:gittree-lifecycle} shows the full lifecycle state machine. The
public CLI presents this lifecycle through \texttt{initialise}, \texttt{pull},
and the READY-state sync commands; lower-level load/expand/validate stages are
internal steps of those workflows.

\begin{center}
\begin{tabular}{lll}
\textbf{Action} & \textbf{Minimum state} & \textbf{Produces} \\
\hline
\texttt{initialise(.cgs)}              & none        & clone tree + \textbf{READY} \\
\texttt{clean-init(.cgs)}              & none        & purge generated state + clone tree + \textbf{READY} \\
\texttt{purge(.cgs)}                   & none        & removed generated root-level clone state \\
\texttt{initialise(.gts)}              & none        & loaded/restored \textbf{READY} \\
\texttt{pull(.cgs|.gts)}               & READY       & resynchronised \textbf{READY} \\
\texttt{checkout(.gts)}                & READY       & \textbf{READY} \\
\texttt{add}                           & \textbf{READY only} & READY \\
\texttt{commit <msg>}                  & \textbf{READY only} & READY \\
\texttt{push}                          & \textbf{READY only} & READY + updated hash \\
\texttt{tag <name>}                    & \textbf{READY only} & READY + updated tag \\
\texttt{freeze}                        & \textbf{READY only} & next \texttt{.gts} id \\
\hline
\end{tabular}
\end{center}

For \texttt{.cgs} refs, repositories can be declared by branch or by tag.
Format validation is static and does not resolve either reference remotely.
The runtime layer selects the declared target (tag takes precedence when both
are present) and checks its availability through \texttt{GitRunner}.

The authoring boundary is explicitly
\texttt{PARSE -> NORMALIZE -> VALIDATE}.  TOML parsing produces authoring data;
normalization expands \texttt{provider:owner/repository} identifiers and fills
\texttt{main} branch defaults, \texttt{ssh}, automatic nested discovery, and
relative paths.  Validation receives only the complete canonical
\texttt{CgsDocument}.  Authoring syntax is therefore not the internal
representation.

File and CLI authoring converge at this boundary. The CLI collects
\texttt{--project} and repeatable \texttt{--repo} values, then calls the
non-interactive \texttt{ComplexGitSyncClient.configure()} Python API. Python
callers can use the same method directly. The facade delegates to
\texttt{CgsDocument}; neither it nor the CLI owns repository grammar or
normalization. Loading an equivalent \texttt{.cgs} file and using either
definition path produce semantically identical canonical documents.
\texttt{create-cgs} runs the same offline pipeline and serializes the result
without entering Git runtime.

The boundary is also responsible for the reverse projection.
\texttt{CgsDocument.to\_git\_tree()} creates the reference model and
\texttt{CgsDocument.from\_git\_tree()} converts a reference or working tree
back to canonical configuration. \texttt{GitTree.to\_cgs()} remains only as a
convenience delegate; it does not assemble repository tables or format TOML.
The minimal serializer then removes deterministically reconstructible defaults.
After the generated TOML is parsed and normalized again, its canonical
representation must equal the one before serialization. Byte equality is not
part of this contract.

Repository placement is parent-relative. At the top level the parent path is
the project root; during nested discovery it is the repository whose
\texttt{.cgs} is being expanded. Thus \texttt{../sibling} resolves to an
existing sibling directory. Discovery compares normalized absolute paths and
keeps an existing registry entry instead of inserting a duplicate.

Placement and traversal are independent. \texttt{nested\_config = auto}
searches the repository root for exactly one \texttt{*.cgs},
\texttt{disabled} stops discovery through that reference, and a relative
\texttt{.cgs} path selects an exact file inside the repository. Explicit paths
may not escape the repository root.

The textual \texttt{provider:owner/repository} grammar belongs exclusively to
\texttt{cgs\_format.py}. Its \texttt{parse\_repo\_id()} function owns syntax validation
and splitting. Normalization uses that function, and CLI repository arguments
must reuse it rather than introduce another parser.

The complete \texttt{.cgs} format pipeline is deterministic and offline.
Parsing, normalization, validation, and serialization do not probe repositories
or execute Git. Remote existence and branch/tag availability are resolved only
after the validated document enters the runtime layer, through explicit
\texttt{GitRunner} operations.

Provider behavior remains in \texttt{git\_repo.py}: the
\texttt{GitProvider} enum, canonical provider set, known-host map, and
\texttt{RepoAddress} URL composer operate on already-separated identity
fields. GitHub, GitLab, and Codeberg map respectively to \texttt{github.com},
\texttt{gitlab.com}, and \texttt{codeberg.org}. \texttt{custom} has no
known-host entry and requires an explicit \texttt{gitprovider\_url}; no
fallback host is guessed.

Before any mutating workflow step (\texttt{commit}, \texttt{push},
\texttt{tag}, or \texttt{freeze}), the action layer now runs a workspace
preflight validator. It checks dirty worktrees, detached HEADs, missing
remotes, branch divergence, unresolved merges, and stale recorded
\texttt{commit\_sha} values. Diagnostics are severity-based:
warnings report actionable-but-allowed states, while blocking errors stop the
operation. \texttt{tag} still requires a clean worktree and a non-existing tag.

\subsection*{Formal \texttt{.gts} snapshot specification (T31)}

\texttt{.gts} now carries explicit schema and deterministic hash metadata in the
\texttt{[document]} table:
\texttt{schema\_version = "1.1"},
\texttt{hash\_algorithm = "sha256"}, and
\texttt{snapshot\_hash = "<64-hex>"}.

The canonical digest is computed from a deterministic payload made of
\texttt{project}, \texttt{tree\_state}, and sorted \texttt{repo\_state} records,
excluding volatile generation metadata (\texttt{generated\_at} and
\texttt{command\_origin}). Non-root entries require
\texttt{parent\_absolute\_path}; READY repositories require \texttt{commit\_sha}.
This makes \texttt{.gts} a stable checkpoint primitive for repeatable workspace
state comparison.

Figure~\ref{fig:ops-sequence} shows the internal sequence of the four
tree-wide synchronisation actions implemented in \texttt{operations.py}.

\begin{figure}[ht]
\centering
\caption{Synchronisation operations: \texttt{pull}/\texttt{checkout\_tree}
  process \texttt{ROOT -> PARENT -> LEAF}; \texttt{add\_tree},
  \texttt{commit\_tree}, and \texttt{push\_tree} iterate
  \texttt{LEAF -> PARENT -> ROOT} so that inner repos are synchronized before
  their parents.}
\label{fig:ops-sequence}
% operations_sequence.tex — TikZ figure showing the synchronized add/checkout/commit/push flow,
% with a footer note for the freeze + tag snapshot workflow.
% Included by docs/architecture.tex (Figure ops-sequence).

\begin{tikzpicture}[
  opbox/.style={
    rectangle, rounded corners=4pt,
    minimum width=3.6cm, minimum height=1.0cm,
    draw, thick, align=center,
    font=\ttfamily\footnotesize,
    inner sep=5pt,
  },
  operationstep/.style={opbox, fill=green!10, draw=green!50},
  gatebox/.style={opbox, fill=yellow!15, draw=yellow!60!black,
                  minimum width=2.4cm, minimum height=0.8cm,
                  font=\sffamily\scriptsize},
  orderbox/.style={opbox, fill=blue!8, draw=blue!40,
                   minimum width=2.6cm, minimum height=0.8cm},
  arr/.style={-{Stealth[length=6pt]}, thick},
  sidarr/.style={-{Stealth[length=5pt]}, thick, gray!70},
  seclabel/.style={font=\sffamily\bfseries\small},
]

% ---- checkout_tree column ------------------------------------------------

\node[seclabel] (co-title) at (0, 0) {\texttt{checkout\_tree}};

\node[gatebox] (co-gate) at (0,-1.0) {assert READY};

\node[operationstep] (co-propagate) at (0,-2.4) {
  \textbf{propagate\_global\_branch}\\
  set target ref on\\every entry (in-memory)
};

\node[operationstep] (co-create) at (0,-4.1) {
  \textbf{create\_global\_branch}\\
  git branch if missing\\(parent-first)
};

\node[operationstep] (co-checkout) at (0,-5.8) {
  \textbf{git checkout}\\
  each repo, parent-first\\refresh entry state
};

\node[orderbox] (co-order) at (0,-7.2) {root → parent → leaf};

\draw[arr] (co-title.south) -- (co-gate.north);
\draw[arr] (co-gate.south)  -- (co-propagate.north);
\draw[arr] (co-propagate.south) -- (co-create.north);
\draw[arr] (co-create.south)    -- (co-checkout.north);
\draw[arr] (co-checkout.south)  -- (co-order.north);

% ---- add_tree column -----------------------------------------------------

\node[seclabel] (ad-title) at (4.3, 0) {\texttt{add\_tree}};

\node[gatebox] (ad-gate) at (4.3,-1.0) {assert READY};

\node[operationstep] (ad-stage) at (4.3,-3.25) {
  \textbf{git add --all}\\
  each repo\\(leaf-first)
};

\node[orderbox] (ad-order) at (4.3,-7.2) {leaf → parent → root};

\draw[arr] (ad-title.south) -- (ad-gate.north);
\draw[arr] (ad-gate.south)  -- (ad-stage.north);
\draw[arr] (ad-stage.south) -- (ad-order.north);

% ---- commit_tree column --------------------------------------------------

\node[seclabel] (cm-title) at (8.6, 0) {\texttt{commit\_tree}};

\node[gatebox] (cm-gate) at (8.6,-1.0) {assert READY};

\node[operationstep] (cm-stage) at (8.6,-2.4) {
  \textbf{git add --all}\\(when stage\_all=True)
};

\node[operationstep] (cm-skip) at (8.6,-4.1) {
  skip if no staged\\changes
};

\node[operationstep] (cm-commit) at (8.6,-5.8) {
  \textbf{git commit}\\
  each repo with staged\\changes; refresh SHA
};

\node[orderbox] (cm-order) at (8.6,-7.2) {leaf → parent → root};

\draw[arr] (cm-title.south) -- (cm-gate.north);
\draw[arr] (cm-gate.south)  -- (cm-stage.north);
\draw[arr] (cm-stage.south) -- (cm-skip.north);
\draw[arr] (cm-skip.south)  -- (cm-commit.north);
\draw[arr] (cm-commit.south)-- (cm-order.north);

% ---- push_tree column ----------------------------------------------------

\node[seclabel] (pu-title) at (12.9, 0) {\texttt{push\_tree}};

\node[gatebox] (pu-gate) at (12.9,-1.0) {assert READY};

\node[operationstep] (pu-push) at (12.9,-3.25) {
  \textbf{git push}\\
  remote = entry.remote\_name\\
  (default: "origin")\\
  branch = resolved\_ref\_name\\
  refresh hash in state
};

\node[orderbox] (pu-order) at (12.9,-7.2) {leaf → parent → root};

\draw[arr] (pu-title.south) -- (pu-gate.north);
\draw[arr] (pu-gate.south)  -- (pu-push.north);
\draw[arr] (pu-push.south)  -- (pu-order.north);

% ---- shared "tree stays READY" footer -----------------------------------

\node[gatebox, minimum width=8.0cm]
    (footer) at (6.45,-8.2) {tree remains \textbf{READY} on success};

\node[orderbox, minimum width=5.0cm]
    (freeze-note) at (12.3,-8.2) {\texttt{commit}/\texttt{push}/\texttt{freeze} preflight: detached HEAD, remotes, divergence, merges, submodules, snapshot SHA drift;\\\texttt{freeze} writes \texttt{freeze\_manifest} invariants + ledger checkpoint and requires an absent release tag;\\CLI \texttt{--dry-run} on \texttt{add|commit|push|freeze} prints \texttt{plan\_actions}/\texttt{plan\_order} without mutation.};

\draw[sidarr] (co-order.south) -- (footer.north west);
\draw[sidarr] (ad-order.south) -- (footer.north);
\draw[sidarr] (cm-order.south) -- (footer.north east);
\draw[sidarr] (pu-order.south) -- (freeze-note.north);

\end{tikzpicture}

\end{figure}

% -----------------------------------------------------------------------
\section{Tier 3 — Client / API}
\label{sec:tier-client}

\texttt{ComplexGitSyncClient} is the single public facade.  It:

\begin{itemize}
  \item holds references to \texttt{Orchestre}, \texttt{GitRunner},
    \texttt{RuntimeStateStore}, and the live \texttt{WorkingGitTree};
  \item computes the current \texttt{TreeLifecycleState} on demand and gates
    every action against it;
  \item emits structured log events for every state transition, action start/end,
    fallback decision, and \texttt{.gts} write/load.
\end{itemize}

\texttt{Orchestre} is the coordination layer between the Client and the tree
models. The CLI (\texttt{cli.py}) collects terminal arguments or interactive
prompt values, then delegates to the reusable Client API. The Client delegates
format semantics to \texttt{cgs\_format.py}; runtime operations remain
separate.

Figure~\ref{fig:client-gating} shows how the Client gates action access based
on the live tree state.

\begin{figure}[ht]
\centering
\caption{Client state-gating: the Client reads \texttt{TreeLifecycleState}
  before delegating to an action.  READY-gated actions (\texttt{checkout},
  \texttt{add}, \texttt{commit}, \texttt{push}, \texttt{tag},
  \texttt{freeze}) are blocked unless the tree is \texttt{READY}.}
\label{fig:client-gating}
% client_gating.tex — TikZ figure showing Client state-gating of actions
% Included by docs/architecture.tex (Figure client-gating).

\begin{tikzpicture}[
  box/.style={
    rectangle, rounded corners=4pt,
    minimum width=3.8cm, minimum height=1.0cm,
    draw, thick, align=center,
    font=\ttfamily\footnotesize,
    inner sep=5pt,
  },
  clientbox/.style={box, fill=blue!12, draw=blue!40},
  statebox/.style={box, fill=yellow!20, draw=yellow!60!black, minimum width=3.2cm},
  actionbox/.style={box, fill=green!12, draw=green!40},
  blockedbox/.style={box, fill=red!10, draw=red!40, minimum width=3.2cm},
  gatebox/.style={
    diamond, draw=gray!60, thick, fill=gray!8,
    minimum width=2.8cm, minimum height=1.4cm,
    font=\sffamily\scriptsize, align=center,
    inner sep=2pt,
  },
  arr/.style={-{Stealth[length=6pt]}, thick},
  redarr/.style={-{Stealth[length=6pt]}, thick, red!60},
  greenarr/.style={-{Stealth[length=6pt]}, thick, green!50!black},
]

% Client
\node[clientbox] (client) at (0,0) {
  \textbf{ComplexGitSyncClient}\\
  \textit{caller invokes action}
};

% Diamond gate
\node[gatebox] (gate) at (0,-2.0) {
  check\\
  TreeLifecycle\\
  State
};

% READY state label
\node[statebox] (readystate) at (3.5,-2.0) {
  TreeLifecycle\\State = \textbf{READY}
};

% READY-gated actions
\node[actionbox] (mutation) at (6.5,-2.0) {
  \textbf{READY-gated actions}\\
  checkout · commit · push\\
  tag · freeze\_release
};

% Any-state actions
\node[actionbox] (bootstrap) at (0,-4.0) {
  \textbf{Bootstrap actions}\\
  clone · restart\\
  launch\_release\\
  \textit{→ must reach READY}
};

% Non-READY blocked
\node[blockedbox] (blocked) at (-4.0,-2.0) {
  \textbf{BLOCKED}\\
  log gating-refusal event\\
  raise GitSyncError
};

% Arrows
\draw[arr] (client.south) -- (gate.north);
\draw[greenarr] (gate.east) -- node[above, font=\scriptsize]{READY} (readystate.west);
\draw[arr] (readystate.east) -- (mutation.west);
\draw[arr] (gate.south) -- node[right, font=\scriptsize]{any state} (bootstrap.north);
\draw[redarr] (gate.west) -- node[above, font=\scriptsize]{not READY} (blocked.east);

\end{tikzpicture}

\end{figure}

% -----------------------------------------------------------------------
\section{Module Layout by Tier}
\label{sec:module-layout}

\begin{lstlisting}[language=bash,caption={Source modules grouped by architectural tier}]
src/ComplexGitSync/
  # --- Tier 1: Core Data (centred on GitTree / GitRepo) ---
  git_repo.py                  GitRepo, WorkingRepo, RepoAddress,
                               RepoNode, enums
  git_tree.py                  GitTree, WorkingGitTree,
                               ProjectTreeState, tree utilities,
                               sync_gitignore (.gitignore maintenance)

  # --- Tier 2: Actions ---
  operations.py                propagate_global_branch, create_global_branch
                               checkout_tree, commit_tree, push_tree
  orchestre.py                 discover_nested_configs(),
                               build_registry_from_cgs_document, render helpers

  # --- Tier 3: Client / API ---
  orchestre.py                 Orchestre, ComplexGitSyncClient
  cli.py                       build_parser / main
  orchestre.py                 GitRunner, RuntimeStateStore,
                               CommandRunLogger + create_run_logger
  master.py                    MasterConfig (workspace-local Git identity
                               for ComplexGitSync-authored commits)

  # --- Document layer (cross-cutting) ---
  config_document.py           ConfigDocument (format-neutral base)
  cgs_format.py                CgsDocument (.cgs parse/normalize/validate/serialize)
  orchestre.py                 GtsDocument
  errors.py                    exception hierarchy
\end{lstlisting}

% -----------------------------------------------------------------------
\section{Document Hierarchy}
\label{sec:document-hierarchy}

Figure~\ref{fig:document-hierarchy} shows the document class hierarchy.

All persistent documents share the \texttt{ConfigDocument} base class, which
provides format-agnostic serialisation (\texttt{to\_toml}, \texttt{to\_json},
\texttt{to\_yaml}) and factory methods (\texttt{from\_toml}, \texttt{from\_json},
\texttt{from\_yaml}).  Subclasses add document-specific validation and
convenience properties.  The base is defined in
\texttt{config\_document.py}; \texttt{cgs\_format.py} owns the complete
\texttt{.cgs}-specific boundary, while \texttt{orchestre.py} consumes validated
documents and retains runtime/orchestration responsibilities.

\begin{itemize}
  \item \textbf{CgsDocument} (\texttt{.cgs}) — the authoring spec; describes
    the static project topology.  \emph{Never} a runtime snapshot.
  \item \textbf{GtsDocument} (\texttt{.gts}) — a generated snapshot capturing
    the exact state of the tree including commit SHAs and absolute paths.
    \emph{Never} hand-edited.
  \item \textbf{Local Git Register} (\texttt{.lgr}) — the project-local
    register that assigns one stable local id to each generated \texttt{.gts}
    and tracks the current snapshot pointer.
\end{itemize}

% -----------------------------------------------------------------------
\section{Registry Model}
\label{sec:registry-model}

Figure~\ref{fig:registry-model} shows the registry model.

\texttt{WorkingGitTree} is the authoritative in-memory graph.  It maps
\emph{repo IDs} (colon-separated path strings such as
\texttt{root:deps/child-repo}) to \texttt{WorkingRepo} instances.
\texttt{WorkingRepo} holds both the static identity declared in
\texttt{.cgs} and the dynamic state updated during operations.

Builder functions in \texttt{orchestre.py} translate document objects into a
live working tree:

\begin{lstlisting}[language=Python]
# Load from authoring spec
registry = build_registry_from_cgs_document(document, config_path)

# Load from snapshot
registry = build_registry_from_gts_document(snapshot_document)

# Serialise to snapshot
gts_doc = build_gts_document_from_registry(
    registry, command_origin="clone", source_cgs_path=path
)
\end{lstlisting}

% -----------------------------------------------------------------------
% Include all TikZ figures defined in figures/class_diagram.tex
% -----------------------------------------------------------------------

Figure~\ref{fig:core-model} is deliberately kept at Tier granularity rather
than diagramming all of \texttt{src/ComplexGitSync/}'s current modules —
see \texttt{docs/DevGuide/architecture.md} for the dev-facing companion
that carries the full per-module Ring-level breakdown and dependency
graph this chapter does not attempt to reproduce.

% class_diagram.tex  — TikZ source for the ComplexGitSync class overview figure.
%
% This file is \input{} from docs/architecture.tex.  It draws four related
% diagrams on separate pages so each fits comfortably within an A4 margin:
%
%   Figure 1 — Class overview by tier (Tier 1/2/3 headline classes, plus the
%              Document-layer and cli/-adapter modules that sit outside the
%              Tier model — every src/ComplexGitSync/ module is named
%              somewhere in this figure, see docs/DevGuide/architecture.md
%              §2 for the full per-module dependency graph)
%   Figure 2 — GitTree lifecycle state machine
%   Figure 3 — Document hierarchy (ConfigDocument + the ConfigDocumentIOMixin
%              file-I/O split)
%   Figure 3 — Registry model (WorkingGitTree and supporting types)

% ============================================================
%  Shared TikZ style definitions
% ============================================================

\usetikzlibrary{calc}

\tikzset{
  %— UML-style class box with three compartments
  classbox/.style={
    rectangle split,
    rectangle split parts=3,
    rectangle split part fill={blue!10,white,white},
    draw=blue!40,
    thick,
    font=\ttfamily\footnotesize,
    text width=4.8cm,
    align=left,
    inner sep=4pt,
  },
  enumbox/.style={
    rectangle split,
    rectangle split parts=2,
    rectangle split part fill={orange!15,white},
    draw=orange!50,
    thick,
    font=\ttfamily\footnotesize,
    text width=3.8cm,
    align=left,
    inner sep=4pt,
  },
  %— Arrows
  hasarrow/.style={-{Stealth[length=6pt]}, thick},
  inharrow/.style={-{Triangle[open,length=7pt]}, thick},
  comporr/.style={-{Diamond[length=7pt]}, thick},
}

% ============================================================
%  Figure 1 — Class overview by tier
% ============================================================
%
% Same three-tier grouping as figures/three_tier.tex (blue = Tier 3, green =
% Tier 2, orange = Tier 1), but drawn at class granularity: one or two
% headline classbox(es) per tier plus a compact list of every other module
% that tier owns.  The two module groups the book's Tier model never
% described — the cross-cutting Document layer and the Ring-4 cli/ adapter
% package (see docs/DevGuide/architecture.md §1's Tier<->Ring mapping) — get
% their own smaller bands below, so every one of the ~27 modules in
% src/ComplexGitSync/ is named once, either as a classbox or in a module
% list.  Full per-module contracts and the real import graph live in
% docs/DevGuide/architecture.md §2-§3; this figure stays at the zoom level a
% book chapter read start-to-finish needs.

\begin{figure}[p]
\centering
\caption{Class overview by tier.  Headline classes for each of the three
  tiers, plus the cross-cutting Document layer and the \texttt{cli/}
  adapter package that the Tier model never described.  Every module in
  \texttt{src/ComplexGitSync/} is named here, either as a class box or in a
  tier's module list; see \texttt{docs/DevGuide/architecture.md} for the
  full current module-by-module Ring-level dependency graph.}
\label{fig:core-model}

\begin{tikzpicture}[
  font=\sffamily,
  tiertitle/.style={font=\sffamily\bfseries\small},
  modlist/.style={font=\ttfamily\scriptsize, align=left, text width=9.6cm},
  smallnote/.style={font=\ttfamily\scriptsize, align=left, text width=6.4cm},
]

% -------- Tier 3 — Client / API --------
\node[tiertitle] (t3title) at (0,0) {Tier 3 --- Client / API};
\node[classbox, below=0.18cm of t3title.south west, anchor=north west] (client) {
  \textbf{ComplexGitSyncClient}
  \nodepart{two}
  orchestre: Orchestre\\
  git\_runner: GitRunner\\
  state\_store: RuntimeStateStore\\
  registry: WorkingGitTree | None\\
  run\_logger: CommandRunLogger | None
  \nodepart{three}
  configure() / initialise()\\
  pull() / clone() / checkout()\\
  add() / git(tree, cmd)\\
  freeze() / verify()
};
\node[classbox, right=0.5cm of client] (orchestreN) {
  \textbf{Orchestre}
  \nodepart{two}
  git\_tree: GitTree
  \nodepart{three}
  register\_repo(repo)
};
\node[modlist, below=0.3cm of client.south west, anchor=north west] (t3mods) {
  also in \texttt{orchestre.py}: \textit{CommandRunLogger} \textperiodcentered{}
  \textit{RuntimeStateStore} \textperiodcentered{} \textit{LocalGitRegister}
  (\texttt{.lgr} writer) \textperiodcentered{} \textit{SyncLedger}
};

\begin{pgfonlayer}{background}
\node[fill=blue!12, draw=blue!40, thick, rounded corners=6pt, inner sep=0.18cm,
      fit=(t3title)(client)(orchestreN)(t3mods)] (tier3band) {};
\end{pgfonlayer}

% -------- Tier 2 — Actions --------
\node[tiertitle, below=0.5cm of tier3band.south west, anchor=north west] (t2title) {Tier 2 --- Actions};
\node[classbox, below=0.18cm of t2title.south west, anchor=north west] (runner) {
  \textbf{GitRunner}
  \nodepart{two}
  executable: str = ``git''
  \nodepart{three}
  clone() / checkout()\\
  commit() / push() / pull()\\
  create\_branch() / create\_tag()\\
  status\_porcelain() / branch\_tracking\_state()
};
\node[modlist, right=0.6cm of runner] (t2mods) {
  supporting modules (\texttt{Tier 2 --- Actions}):\\
  \textperiodcentered{} \texttt{operations.py} --- leaf/parent-first Git ops\\
  \textperiodcentered{} \texttt{registry.py} --- build/serialize \texttt{WorkingGitTree}\\
  \textperiodcentered{} \texttt{paths.py}, \texttt{discovery.py}\\
  \textperiodcentered{} \texttt{state\_store.py}, \texttt{snapshot\_resolver.py}\\
  \textperiodcentered{} \texttt{ledger\_store.py}
};

\begin{pgfonlayer}{background}
\node[fill=green!10, draw=green!40, thick, rounded corners=6pt, inner sep=0.18cm,
      fit=(t2title)(runner)(t2mods)] (tier2band) {};
\end{pgfonlayer}

\draw[-{Stealth[length=8pt]}, thick, gray]
  (tier3band.south) -- node[right, font=\scriptsize\itshape]{calls (state-gated)} (tier2band.north);

% -------- Tier 1 — Core Data --------
\node[tiertitle, below=0.5cm of tier2band.south west, anchor=north west] (t1title) {Tier 1 --- Core Data};
\node[classbox, below=0.18cm of t1title.south west, anchor=north west] (gittree) {
  \textbf{GitTree} / \textbf{WorkingGitTree}
  \nodepart{two}
  repos: dict[str, GitRepo]\\
  \textit{(WorkingGitTree: dict[str, WorkingRepo],}\\
  \textit{ + lifecycle\_state: TreeLifecycleState)}
  \nodepart{three}
  add\_repo() / force\_repo\_sha()\\
  recompute\_tree\_state()\\
  \textit{WorkingGitTree(GitTree)}
};
\node[classbox, right=0.5cm of gittree] (gitrepo) {
  \textbf{GitRepo} / \textbf{WorkingRepo}
  \nodepart{two}
  project\_name / project\_owner\_name\\
  gitprovider / access\_protocol\\
  commit\_sha
  \nodepart{three}
  resolved\_group\_name\\
  \textit{WorkingRepo(GitRepo)}: repo\_id,\\
  \textit{node\_type, repo\_lifecycle\_state, ...}
};
\node[modlist, below=0.3cm of gittree.south west, anchor=north west] (t1mods) {
  also Tier 1: \texttt{master.py}'s \textit{MasterConfig} (workspace-local
  commit identity) \textperiodcentered{} \textit{RepoAddress} (URL composer,
  \texttt{git\_repo.py})
};

\begin{pgfonlayer}{background}
\node[fill=orange!10, draw=orange!40, thick, rounded corners=6pt, inner sep=0.18cm,
      fit=(t1title)(gittree)(gitrepo)(t1mods)] (tier1band) {};
\end{pgfonlayer}

\draw[-{Stealth[length=8pt]}, thick, gray]
  (tier2band.south) -- node[right, font=\scriptsize\itshape]{reads / mutates} (tier1band.north);

% -------- Cross-cutting: Document layer + cli/ adapter --------
\node[tiertitle, below=0.5cm of tier1band.south west, anchor=north west] (docttitle) {Document layer (cross-cutting)};
\node[smallnote, below=0.18cm of docttitle.south west, anchor=north west] (docmods) {
  \textbf{ConfigDocument} (base) $\to$ \textbf{CgsDocument} / \textbf{GtsDocument},\\
  each also mixing in \textit{ConfigDocumentIOMixin}\\
  (\texttt{config\_document\_io.py}) for \texttt{from/to\_toml|json|yaml}.\\
  Plus \texttt{errors.py}, and the no-Tier-home register/status\\
  internals: \texttt{integrity.py}, \texttt{ledger\_entry.py},\\
  \texttt{status\_render.py}
};

\begin{pgfonlayer}{background}
\node[fill=gray!12, draw=gray!50, thick, rounded corners=6pt, inner sep=0.18cm,
      fit=(docttitle)(docmods)] (docband) {};
\end{pgfonlayer}

\node[tiertitle, right=0.6cm of docband, anchor=north west] (clititle) at (docband.north east) {\texttt{cli/} adapter (Ring 4, outside the Tier model)};
\node[smallnote, below=0.18cm of clititle.south west, anchor=north west] (climods) {
  \texttt{cli/\_\_init\_\_.py} (parser + dispatch)\\
  \texttt{cli/minimalist.py}, \texttt{cli/configuration.py},\\
  \texttt{cli/expert.py} (one module per command group)\\
  \texttt{cli/\_shared.py} (cross-command helpers)\\[2pt]
  \textit{argument/prompt collection only --- calls into}\\
  \textit{Tier 3 above, never touches Git directly}
};

\begin{pgfonlayer}{background}
\node[fill=red!8, draw=red!40, thick, dashed, rounded corners=6pt, inner sep=0.18cm,
      fit=(clititle)(climods)] (cliband) {};
\end{pgfonlayer}

\end{tikzpicture}

\noindent\footnotesize\textit{Not shown: \texttt{\_\_init\_\_.py} (package
re-exports) and \texttt{\_\_main\_\_.py} (entry-point stub) --- package
assembly, not part of any tier or ring.}
\end{figure}

% ============================================================
%  Figure 2 — GitTree lifecycle states
% ============================================================

\begin{figure}[ht]
\centering
\caption{GitTree lifecycle state machine.  The simplified user-facing lifecycle
  is \texttt{initialise(.cgs)} $\rightarrow$ clone $\rightarrow$ \texttt{READY};
  \texttt{initialise(.gts)} $\rightarrow$ restore $\rightarrow$ \texttt{READY};
  \texttt{pull(.cgs|.gts)} $\rightarrow$ resync $\rightarrow$ \texttt{READY}.
  Synchronized Git operations (\texttt{checkout},
  \texttt{add}, \texttt{git("commit")}, \texttt{git("push")},
  \texttt{git("tag")},
  \texttt{freeze}) are gated on \texttt{READY}.}
\label{fig:gittree-lifecycle}

\begin{tikzpicture}[
  node distance=1.4cm and 2.0cm,
  state/.style={rectangle, rounded corners=4pt, draw=blue!50, fill=blue!8,
                font=\ttfamily\small, minimum width=2.6cm, minimum height=0.8cm,
                inner sep=4pt},
  sidestate/.style={rectangle, rounded corners=4pt, draw=orange!60, fill=orange!10,
                    font=\ttfamily\small, minimum width=2.6cm, minimum height=0.8cm,
                    inner sep=4pt},
  trans/.style={-{Stealth[length=6pt]}, thick},
]

\node[state] (unloaded) {UNLOADED};
\node[state, right=of unloaded] (declared)  {DECLARED};
\node[state, right=of declared] (pending)   {PENDING};
\node[state, right=of pending]  (ready)     {READY};

\node[sidestate, below=1.2cm of declared] (partial) {PARTIAL};
\node[sidestate, below=1.2cm of pending]  (error)   {ERROR};

\draw[trans] (unloaded) -- node[above, font=\scriptsize]{load(.cgs)} (declared);
\draw[trans] (declared) -- node[above, font=\scriptsize]{discover nested} (pending);
\draw[trans] (pending)  -- node[above, font=\scriptsize]{clone/validate} (ready);
\draw[trans, dashed] (declared.south) -- (partial.north);
\draw[trans, dashed] (pending.south)  -- (error.north);
\draw[trans, bend right=20] (ready.north) to node[above, font=\scriptsize]{freeze / next .gts id + .lgr update} (declared.north);

\end{tikzpicture}
\end{figure}

% ============================================================
%  Figure 3 — Document class hierarchy
% ============================================================

\begin{figure}[ht]
\centering
\caption{Document class hierarchy.  \texttt{ConfigDocument} (\texttt{config\_document.py})
  is the abstract, I/O-free base; \texttt{ConfigDocumentIOMixin}
  (\texttt{config\_document\_io.py}) supplies the file-based
  \texttt{from/to\_toml|json|yaml} methods.  Each concrete document type
  combines both via ordinary multiple inheritance.}
\label{fig:document-hierarchy}

\begin{tikzpicture}[node distance=1.2cm and 1.8cm]

\node[classbox] (base) {
  \textbf{ConfigDocument}
  \nodepart{two}
  config\_document.py\\
  \_data: dict
  \nodepart{three}
  read(key) / get(key)\\
  validate()\\
  to\_dict() / from\_dict(data)
};

\node[classbox, right=2.4cm of base] (iomixin) {
  \textbf{ConfigDocumentIOMixin}
  \nodepart{two}
  config\_document\_io.py
  \nodepart{three}
  from\_toml / from\_json / from\_yaml\\
  to\_toml / to\_json / to\_yaml\\
  print()
};

\node[classbox, below left=1.6cm and 0.6cm of base] (cgs) {
  \textbf{CgsDocument}
  \nodepart{two}
  cgs\_format.py\\
  DOCUMENT\_KIND = ``cgs''
  \nodepart{three}
  project\_name\\
  default\_branch\\
  repos\\
  normalize authoring syntax\\
  to\_git\_tree() / from\_git\_tree()\\
  to\_authoring\_dict() / to\_toml()\\
  runtime\_setting(key)
};

\node[classbox, below=1.6cm of base] (gts) {
  \textbf{GtsDocument}
  \nodepart{two}
  gts\_document.py\\
  DOCUMENT\_KIND = ``gts''
  \nodepart{three}
  schema\_version\\
  snapshot\_hash\\
  compute\_snapshot\_hash()\\
  ensure\_snapshot\_hash()\\
  lifecycle\_state\\
  is\_ready\\
  repo\_states
};

\draw[inharrow] (cgs.north) -- (base.south west);
\draw[inharrow] (gts.north) -- (base.south);
\draw[inharrow, dashed] (cgs.north east) -- ($(iomixin.south)+(-0.6,0)$);
\draw[inharrow, dashed] (gts.east) -- (iomixin.south);

\end{tikzpicture}
\end{figure}

% ============================================================
%  Figure 4 — Registry model
% ============================================================

\begin{figure}[ht]
\centering
\caption{Registry model.  \texttt{WorkingGitTree} is the authoritative
  in-memory graph; \texttt{WorkingRepo} \emph{extends} \texttt{GitRepo}
  (Figure~\ref{fig:core-model}) rather than duplicating its identity
  fields, adding only the mutable tree-position and synchronisation state
  below.}
\label{fig:registry-model}

\begin{tikzpicture}[node distance=1.2cm and 2.0cm]

% WorkingGitTree
\node[classbox] (registry) {
  \textbf{WorkingGitTree}
  \nodepart{two}
  repos: dict[str, WorkingRepo]\\
  lifecycle\_state: TreeLifecycleState
  \nodepart{three}
  add(entry) / get(id) / values()\\
  children\_of(parent\_id)\\
  is\_complete() / is\_ready()\\
  recompute\_tree\_state()
};

% WorkingRepo
\node[classbox, right=2.2cm of registry] (entry) {
  \textbf{WorkingRepo} \textit{(extends GitRepo)}
  \nodepart{two}
  repo\_id / name / node\_type\\
  parent\_id / absolute\_path\\
  current\_ref\_kind / \_name\\
  target\_ref\_kind / \_name\\
  resolved\_ref\_kind / \_name\\
  repo\_lifecycle\_state\\
  sync\_state / discovery\_state\\
  \textit{(+ inherited: project\_name,}\\
  \textit{gitprovider, commit\_sha, ...)}
  \nodepart{three}
  (mutable dataclass)
};

% ProjectTreeState
\node[classbox, below=1.6cm of registry] (treestate) {
  \textbf{ProjectTreeState}
  \nodepart{two}
  lifecycle\_state: TreeLifecycleState\\
  is\_ready: bool\\
  registry\_complete: bool
  \nodepart{three}
  (frozen dataclass)
};

% Enum: NodeType
\node[enumbox, below left=1.6cm and 0.2cm of entry] (nodetype) {
  \textbf{NodeType} \textit{(enum)}
  \nodepart{two}
  ROOT / PARENT / LEAF
};

% Enum: TreeLifecycleState
\node[enumbox, right=1.8cm of treestate] (tlc) {
  \textbf{TreeLifecycleState}
  \nodepart{two}
  UNLOADED / DECLARED\\
  PENDING / READY\\
  PARTIAL / ERROR
};

% Enum: RepoLifecycleState
\node[enumbox, below=1.2cm of nodetype] (rlc) {
  \textbf{RepoLifecycleState}
  \nodepart{two}
  DECLARED / PENDING\\
  READY / FALLBACK\_READY\\
  MISSING / ERROR
};

% Arrows
\draw[comporr] (registry.east) -- node[above, font=\scriptsize]{0..*} (entry.west);
\draw[hasarrow] (registry.south) -- (treestate.north);
\draw[inharrow, dashed] (entry.north) -- ++(0,1.0) node[above, font=\scriptsize\itshape]{GitRepo (Fig.~\ref{fig:core-model})};
\draw[hasarrow] (entry.south west) -- (nodetype.north east);
\draw[hasarrow] (entry.south) -- (rlc.north);
\draw[hasarrow] (registry.south east) -- (tlc.north west);

\end{tikzpicture}
\end{figure}

